\section{THPT Nguyễn Văn Trỗi}

\debai{
Trong mặt phẳng với hệ trục tọa độ $Oxy$, cho tam giác $ABC$ có phương trình đường thẳng $BC:x-y-4=0$. Các điểm $H(2;0), I(3;0)$ lần lượt là trực tâm và tâm đường tròn ngoại tiếp tam giác. Hãy lập phương trình cạnh $AB$, biết điểm $B$ có hoành độ không lớn hơn $3$.}

\debai{
Giải hệ phương trình $\begin{cases} x^3 -3x^2+2 = \sqrt{y^3+3y^2} \\ \sqrt{x-3} = \sqrt{y-x+2}\end{cases}$
}

\debai{
Cho $a,b,c>0$ thỏa mãn $a+b+c=1$, chứng minh rằng\\ 
\cl{$\dfrac{a}{1+bc}+\dfrac{b}{1+ca}+\dfrac{c}{1+ab}\geq \dfrac{9}{10}$}
}

\section{THPT Chuyên ĐH Vinh}

\debai{
Giải bất phương trình $x^2+5x<4 \left(1+ \sqrt{x^3+2x^2-4x} \right)$}

\debai{
Trong mặt phẳng với hệ trục tọa độ $Oxy$, cho hình chữ nhật $ABCD$ có $\widehat{ACD}=\alpha$ với $\cos \alpha = \dfrac{1}{\sqrt{5}}$, điểm $H$ thỏa mãn điều kiện $\vt{HB}=-2\vt{HC}$, $K$ là giao điểm của hai đường thẳng $AH$ và $BD$. Biết $H\left(\dfrac{1}{3}-\dfrac{4}{3}\right), K(1;0)$ và điểm $B$ có hoành độ dương. Tìm tọa độ các điểm $A,B,C,D$.
}

\debai{
Cho $x,y,z$ là các số thực không âm thỏa mãn 
$0 < (x+y)^2 +(y+z)^2 +(z+x)^2 \leq 2$
Tìm giá trị lớn nhất của biểu thức:
$$P=4^x+4^y+4^z+\ln(x^4+y^4+z^4) - \dfrac{3}{4}(x+y+z)^4$$
}
\section{THPT Thủ Đức (TP Hồ Chí Minh)}
\debai{
Trong mặt phẳng với hệ trục tọa độ $Oxy$, cho tam giác $ABC$ nhọn có $A(-1;4)$, các đường cao $AM,CN$ và trực tâm $H$. Tâm đường tròn ngoại tiếp tam giác $HMN$ là $I(2;0)$. Đường thẳng $BC$ đi qua $P(1;-2)$. Tìm tọa độ các đỉnh của tam giác, biết $B$ thuộc đường thẳng $d: x+2y-2=0$
}
\debai{
Giải hệ phương trình 
$\va{(1-y)\sqrt{x^2+2y^2}=x+2y+3xy \\ \sqrt{y+1}+\sqrt{x^2+2y^2}=2y-x}$
}

\debai{
Cho $x,y,z$ là các số thực dương thỏa mãn \\
\cl{$5(x^2+y^2+z^2)=9(xy+2yz+zx)$}
Tìm giá trị lớn nhất của biểu thức:\\
\cl{$P=\dfrac{x}{y^2+z^2}-\dfrac{1}{(x+y+z)^3}$}
}
\section{THPT Nông Cống 1 (Thanh Hóa) lần 2}
\debai{
Trong mặt phẳng với hệ trục tọa độ $Oxy$, cho đường tròn $(C)$ tâm $I$ có hoành độ dương. $(C)$ đi qua điểm $A(-2;3)$ và tiếp xúc với đường thẳng $(d_1):x+y+4=0$ tại điểm $B$. \\
$(C)$ cắt đường thẳng $(d_2):3x+4y-16=0$ tại $C$ và $D$ sao cho $ABCD$ là hình thang $(AD // BC)$ có hai đường chéo $AC,BD$ vuông góc với nhau. Tìm tọa độ các điểm $B,C,D$.}

\debai{
Giải hệ phương trình 
 $\va{\sqrt{x^2+xy+2y^2}+\sqrt{y^2+xy+2x^2}=2(x+y) \\ (8y-6)\sqrt{x-1}=\left(2+\sqrt{y-2} \right) \left(y + 4\sqrt{x-2}+3 \right)}$
}

\debai{
Cho $x,y$ là các số thực không âm thỏa mãn \\
\cl{$\sqrt{2x^2+3xy+4y^2}+\sqrt{2y^2+3xy+4x^2}-3(x+y)^2 \leq 0$}
Tìm giá trị nhỏ nhất của biểu thức:\\
\cl{$P=2(x^3+y^3)+2(x^2+y^2)-xy+\sqrt{x^2+1}+\sqrt{y^2+1}$}
}

\section{THPT Nguyễn Trung Thiên lần 1}
\debai{
Trong mặt phẳng với hệ trục tọa độ $Oxy$, cho tam giác $ABC$ có trung điểm cạnh $BC$ là $M(3;-1)$. Điểm $E(-1;-3)$ nằm trên đường thẳng $\Delta$ chứa đường cao đi qua đỉnh $B$. Đường thẳng $AC$ đi qua $F(1;3)$. Tìm tọa độ các đỉnh của tam giác $ABC$, biết đường tròn ngoại tiếp tam giác $ABC$ có đường kính $AD$ với $D(4;-2)$.}
\\
\debai{
Giải phương trình 
 $(x+2)\left(\sqrt{x^2+4x+7}+1 \right) + x\left(\sqrt{x^2+3}+1\right)=0$
}

\debai{
Cho $x,y,z$ là các số thực dương thỏa mãn $x^2+y^2+z^2=3$.\\
Tìm giá trị nhỏ nhất của biểu thức:\\
\cl{$P=\dfrac{x}{(y+z)^2}+\dfrac{y}{(x+z)^2}+\dfrac{z}{(x+y)^2}$}
}
\section{THPT Lam Kinh}
\debai{
Trong mặt phẳng với hệ trục tọa độ $Oxy$, cho hình bình hành $ABCD$ có diện tích bằng 4. Biết $A(1;0), B(0;2)$ và tâm $I$ của hình bình hành nằm trên đường thẳng $y=x$. Tìm tọa độ các đỉnh $C$ và $D$.}
\\
\debai{
Giải hệ phương trình 
 $\va{x^2+y^2+xy+1=4y\\y(x+y)^2=2x^2+7y+2}$
}

\debai{
Cho $a,b,c$ là ba cạnh của một tam giác. Chứng minh rằng:\\
\cl{$a\left(\dfrac{1}{3a+b}+\dfrac{1}{3a+c}+\dfrac{1}{2a+b+c}\right)+\dfrac{b}{3a+c}+\dfrac{c}{3a+b}< 2$}
}
\section{THPT Cù Huy Cận (Hà Tĩnh)}
\debai{
Trong mặt phẳng với hệ trục tọa độ $Oxy$, cho hình chữ nhật $ABCD$ có điểm $A$ thuộc đường thẳng $d_1:x-y-4=0$, điểm $C(-7;5)$, $M$ là điểm thuộc $BC$ sao cho $MB=3MC$, đường thẳng đi qua $D$ và $M$ có phương trình $d_2:3x-y+18=0$. Xác định tọa độ đỉnh $A,B$ biết điểm $B$ có tung độ dương.}
\\
\debai{
Giải hệ phương trình 
 $\va{x^2+2x-3=x+3\sqrt{x+y+3}\\6x^2+2xy+2\left(\sqrt{x}-1\right)\left(\sqrt{x}=1\right)=3(x^2-y-4)\sqrt[3]{2x^2+xy+3x+2}
}$}

\debai{
Cho $x,y$ là các số thực thỏa mãn: $4x^2+y^2 \leq 8$. Tìm giá trị lớn nhất, giá trị nhỏ nhất của:
\cl{$P=\dfrac{(2x+6)^2+(y+6)^2+4xy-32}{2x+y+6}$}
}
\section{THPT Đa Phúc (Hà Nội)}
\debai{
Trong mặt phẳng với hệ trục tọa độ $Oxy$, cho hình vuông $ABCD$ có $E,F\left(\dfrac{11}{2};3\right)$ lần lượt là trung điểm $AB,AD$. Gọi $K$ là điểm thuộc cạnh $CD$ sao cho $KD=3KC$. Phương trình đường thẳng $EK:19x-8y-18=0$. Xác định tọa độ đỉnh $C$, biết điểm $E$ có hoành độ nhỏ hơn 3.}

\debai{
Giải hệ phương trình 
 $\va{\sqrt{\dfrac{x^2+y^2}{2}}+\sqrt{\dfrac{x^2+xy+y^2}{3}}=x+y\\x\sqrt{2xy+5x+3}=4xy-5x-3}$}

\debai{
Cho $a,b, c$ là các số thực đôi một phân biệt thỏa mãn: $\va{a+b+c=1\\ab+bc+ca>0}$. Tìm giá trị nhỏ nhất của:
$$P=2\left(\sqrt{\dfrac{2}{(a-b)^2}+\dfrac{2}{(b-c)^2}}+\dfrac{1}{|c-a|} \right)+\dfrac{5}{\sqrt{ab+bc+ca}}$$}

\section{THPT Lạng Giang I (Bắc Giang)}
\debai{
Trong mặt phẳng với hệ trục tọa độ $Oxy$, cho tam giác $ABC$. Gọi $E, F$ lần lượt là chân đường cao hạ từ $B,C$. Tìm tọa độ điểm $A$, biết $E(7;1), F\left( \dfrac{11}{5};\dfrac{13}{5} \right)$, phương trình đường thẳng $BC$ là $x+3y-4=0$ và điểm $B$ có tung độ dương.}
\\
\debai{
Giải hệ phương trình 
 $\va{(xy+3)^2+(x+y)^2=8 \\ \dfrac{x}{x^2+1}+\dfrac{y}{y^2+1}=-\dfrac{1}{4}}$}

\debai{
Cho các số thực dương $x,y,z$. Tìm giá trị lớn nhất của biểu thức:\\
\cl{$P=\dfrac{1}{\sqrt{x^2+y^2+z^2+4}}-\dfrac{8}{(x+2)(y+2)(z+2)}$}
}
\section{THPT Lý Tự Trọng (Khánh Hòa)}
\debai{
Trong mặt phẳng với hệ trục tọa độ $Oxy$, cho hai điểm $A(1;2), B(4;1)$ và đường thẳng $d:3x-4y+5=0$. Viết phương trình đường tròn $(C)$ đi qua $A,B$ và cắt đường thẳng $d$ tại $C,D$ sao cho $CD=6$.}
\\
\debai{
Giải hệ phương trình 
 $\va{x^3-6x^2+13x=y^3+y+10 \\ \sqrt{2x+y+5}-\sqrt{3-x-y}=x^3-3x^2-10y+6}$}

\debai{
Cho các số thực không âm $x,y,z$ thỏa mãn điều kiện: $x+y+z=3$. Tìm giá trị nhỏ nhất của biểu thức:\\
\cl{$P=\dfrac{1}{4+2\ln (1+x)-y}+\dfrac{1}{4+2\ln (1+y)-z}+\dfrac{1}{4+2\ln (1+z)-x}$}
}
\section{THPT Quảng Hà}
\debai{
Trong mặt phẳng với hệ trục tọa độ $Oxy$, cho hình vuông $ABCD$ có đỉnh $A(2;2)$. Biết điểm $M(6;3)$ thuộc cạnh $BC$, điểm $N(4;6)$ thuộc cạnh $CD$. Tìm tọa độ đỉnh $C$.}

\debai{
Giải hệ phương trình 
 $\va{x^4+4y^3+(x^4-1)y+4y^2=1 \\ 8y^3+4\sqrt{x^2+1}=x^2+6y+2}$}

\debai{
Cho các số thực dương $x,y,z$ thỏa mãn điều kiện: $x+y+z\geq 3$. Tìm giá trị nhỏ nhất của biểu thức:\\
\cl{$P=\dfrac{x^2}{yz+\sqrt{8+x^3}}+ \dfrac{y^2}{zx+\sqrt{8+y^3}}+ \dfrac{z^2}{xy+\sqrt{8+z^3}}$}}

\section{THPT Thống nhất}
\debai{Giải hệ phương trình 
 $\va{x-2y-\sqrt{xy}=0\\ \sqrt{x-1}-\sqrt{2y-1}=1}$}

\debai{
Trong mặt phẳng với hệ trục tọa độ $Oxy$, cho các điểm $A(1;0), B(-2;4), C(-1;4),D(3;5)$ và đường thẳng $d:3x-y-5=0$. Tìm tọa độ điểm $M$ trên $d$ sao cho hai tam giác $MAB,MCD$ có diện tích bằng nhau.}

\debai{
Cho các số thực dương $a,b,c$ thỏa mãn điều kiện: $a+b+c=1$. Chứng minh rằng:\\
\cl{$\dfrac{a+b^2}{b+c}+ \dfrac{b+c^2}{c+a}+ \dfrac{c+a^2}{a+b}\geq 2$}}

\section{THPT Hồng Quang (Hải Dương)}
\debai{
Trong mặt phẳng với hệ trục tọa độ $Oxy$, cho tam giác $ABC$ có đường cao và đường trung tuyến kẻ từ đỉnh $A$ lần lượt có phương trình $x-3y=0$ và $x+5y=0$. Đỉnh $C$ nằm trên đường thẳng $\Delta: x+y-2=0$ và có hoành độ dương. Tìm tọa độ các đỉnh của tam giác, biết rằng đường trung tuyến kẻ từ $C$ đi qua điểm $E(-2;6)$}
\\\debai{Giải hệ phương trình 
 $\va{x-\dfrac{1}{(x+1)^2}=\dfrac{y}{x+1}-\dfrac{1+y}{y}\\ \sqrt{8y+9}=(x+1)\sqrt{y}+2}$}
\\

\debai{
Cho các số thực dương $x,y,z$ thỏa mãn điều kiện: $x>y$ và $(x+z)(y+z)=1$. Tìm giá trị nhỏ nhất của biểu thức:\\
\cl{$P = \dfrac{1}{(x-1)^2}+\dfrac{4}{(x+z)^2}+\dfrac{4}{(y+z)^2}$}}
\section{THPT Sông Lô (Vĩnh Phúc)}
\debai{
Trong mặt phẳng với hệ trục tọa độ $Oxy$, cho tam giác $ABC$ cân tại $A$, đường trung tuyến kẻ từ đỉnh $A$ là $d:2x+y-3=0$. Đỉnh $B$ thuộc trục hoành, đỉnh $C$ thuộc trục tung và diện tích tam giác $ABC$ bằng $5$. Tìm tọa độ các đỉnh của tam giác.}

\debai{Giải hệ phương trình 
 $\va{x+y+\sqrt{x^2-y^2}=12\\y\sqrt{x^2-y^2}=12}$}


\debai{
Cho các số thực dương $a,b,c$ thỏa mãn điều kiện: $abc=1$. Chứng minh rằng:\\
\cl{$\dfrac{\sqrt{a}}{2=b\sqrt{a}}+\dfrac{\sqrt{b}}{2+c\sqrt{b}}+\dfrac{\sqrt{c}}{2+a\sqrt{c}}\geq 1$}}

\section{THPT chuyên Hà Tĩnh}
\debai{Trong mặt phẳng hệ tọa độ $Oxy$ cho hình bình hành $ABCD$ có $N$ là trung điểm cạnh $CD$ và đường thẳng $BN$ có phương trình là $13x-10y+13=0$, điểm $M(-1;2)$ thuộc đoạn thẳng $AC$ sao cho $AC=4AM$. Gọi $H$ là điểm đối xứng với $N$ qua $C$. Tìm tọa độ các đỉnh $A,B,C,D$ biết rằng $3AC=2AB$ và điểm $H$ thuộc đường thẳng $\Delta: 	2x-3y=0$}

\debai{Giải hệ phương trình $\va{x^2+(y^2-y-1)\sqrt{x^2+2}-y^3+y+2=0 \\ \sqrt[3]{y^2-3}-\sqrt{xy^2-2x-2}+x=0} (x,y\in \mathbb{R})$}

\debai{Cho $a$ là số thực dương thuộc đoạn $[1;2]$. Chứng minh rằng $$(2^a+3^a+4^a)(6^a+8^a+12^a)<24^{a+1}$$
}
\section{THPT chuyên Nguyễn Huệ lần 3}
\debai{Cho đường tròn $(C)$ có phương trình $x^2+y^2-2x-4y+1=0$ và $P(2;1)$. Một đường thẳng $d$ đi qua $P$ cắt đường tròn tại $A$ và $B$. Tiếp tuyến tại $A$ và $B$ của đường tròn cắt nhau tại $M$. Tìm tọa độ của $M$ biết $M$ thuộc đường tròn  $(C_1):x^2+y^2 -6x-4y+11=0$}

\debai{Giải hệ phương trình $\va{x+y+\sqrt{2y-1}+\sqrt{x-y}=5 \\ y^2+2=xy+y} (x,y\in \mathbb{R})$}

\debai{Với $a,b,c$ là các số thực thỏa mãn $a^2+b^2+c^2=3$. Tìm giá trị lớn nhất của biểu thức $$P=a^4+b^4+c^4+3(ab+bc+ca)$$}

\section{THPT chuyên Hùng Vương}
\debai{Trong mặt phẳng hệ tọa độ $Oxy$, cho tam giác $ABC$ có đường trung tuyến $AM$ và đường cao $AH$ lần lượt có phương trình $13x-6y-2=0$, $x-2y-14=0$. Tìm tọa độ các đỉnh của tam giác $ABC$ biết tâm đường tròn ngoại tiếp của tam giác $ABC$ là $I(-6;0)$.}

\debai{Giải bất phương trình $2x+5\sqrt{x}>11+\dfrac{4}{x-2}$}

\debai{Giả sử $a,b,c$ là các số thực dương thỏa mãn $a+b+c=1$. Tìm giá trị nhỏ nhất của biểu thức
$$P=\dfrac{a^2}{(b+c)^2+5bc}+\dfrac{b^2}{(c+a)^2+5ca}-\dfrac{3}{4}(a+b)^2$$}

\section{Chuyên Nguyễn Huệ (Quảng Nam)}

\debai{Trong mặt phẳng hệ tọa độ $Oxy$ cho điểm $M(1;2)$ và hai đường thẳng $d_1:x-2y+1=0$ và $d_2:2x+y+2=0$. Gọi $A$ là giao điểm của $d_1$ và $d_2$. Viết phương trình đường thẳng $d$ đi qua $M$ và cắt các đường thẳng $d_1,d_2$ lần lượt tại $B,C$ ($B,C$ không trùng $A$) sao cho $\dfrac{1}{AB^2}+\dfrac{1}{AC^2}$ đạt giá trị nhỏ nhất.}

\debai{Giải hệ phương trình $\va{3x^2+12y^2+24xy-9(x+2y)\sqrt{2xy}=0
\\
5x^2-7y^2+xy=15
} (x,y \in \mathbb{R})$}

\debai{Cho $a,b$ là các số thực dương thỏa mãn $a+b+ab=3$. Tìm giá trị nhỏ nhất của biểu thức 
$$P=\dfrac{a^2}{1+b}+\dfrac{b^2}{1+a}+\dfrac{ab}{a+b}$$}

\section{Chuyên Lê Quý Đôn (Bình Định)}
\debai{Trong mặt phẳng hệ tọa độ $Oxy$, cho hình thang cân $ABCD$ ($AD\parallel BC$) có phương trình đường thẳng $AB:x-2y+3=0$ và đường thẳng $AC:y-2=0$. Gọi $I$ là giao điểm của $AC$ và $BD$. Tìm tọa độ các đỉnh của hình thang cân $ABCD$, biết $IB=IA\sqrt{2}$, hoành độ $I$ lớn hơn $-3$ và $M(-1;3)$ nằm trên đường thẳng $BD$.}

\debai{Giải hệ phương trình $\va{(1-y)(x-3y+3)-x^2=\sqrt{(y-1)^3}.\sqrt{x} \\ \sqrt{x^2-y}+2\sqrt[3]{x^3-4}=2(y-2)}(x,y \in \mathbb{R})$}

\debai{Cho $x,y$ là hai số thực dương thỏa mãn $2x+3y\le 7$. Tìm giá trị nhỏ nhất của biểu thức 
$$P=2xy+y+\sqrt{5(x^2+y^2)}-24\sqrt[3]{8(x+y)-(x^2+y^2+3)}$$}

\section{Chuyên ĐH Vinh lần 3}
\debai{Trong mặt phẳng với hệ tọa độ $Oxy$, cho tam giác $ABC$ có trọng tâm $G\left(\dfrac{8}{3};0\right)$ và có đường tròn ngoại tiếp là $(C)$ tâm $I$. Biết rằng các điểm $M(0;4), N(4;1)$ lần lượt đối xứng $I$ qua các đường thẳng $AB,AC$, đường thẳng $BC$ đi qua điểm $K(2;-1)$. Viết phương trình đường tròn $(C)$}

\debai{ Giải bất phương trình: \qquad $3(x^2-1)\sqrt{2x+1}<	 2(x^3-x^2) \qquad(1)$. }

\debai{Giả sử $x,y,z$ là các số thực dương thỏa mãn $x+z\le 2y$ và $x^2+y^2+z^2 =1$. Tìm giá trị lớn nhất của biểu thức
$$P=\dfrac{xy}{1+z^2}+\frac{yz}{1+x^2}-y^3 \left(\frac{1}{x^3} +\frac{1}{z^3}\right)$$
}

\section{Chuyên Hùng Vương (Gia Lai)}
\debai{Trong mặt phẳng tọa độ $Oxy$, cho hình chữ nhật $ABCD$ có diện tích bằng 16, các đường thẳng $AB,BC,CD,DA$ lần lượt đi qua các điểm $M(4;5), N(6;5), P(5;2), Q(2;1)$. Viết phương trình đường thẳng $AB$.}

\debai {Giải hệ phương trình: $\va {  x^3-x^2y-y=2x^2-x+2 &  (1)\\
y^2+4\sqrt{x+2} +\sqrt{16-3y}=2x^2-4x+12 & (2)} (x,y \in \mathbb{R})$}

\debai {Cho $x;y;z $ là 3 số thực thuộc đoạn  $[1;2]$. Tìm giá trị nhỏ nhất của biểu thức:
$$P=\dfrac{(x+y)^2}{2(x+y+z)^2-2(x^2+y^2)-z^2}$$}